%
%   В цьому окремому файлі визначимо власні команди,
%   щоб не засмічувати головний файл
%


%                        Custom colors
% --------------------------------------------------------------
\definecolor{grey50}{HTML}{FAFAFA}
\definecolor{bluegray400}{HTML}{78909C}

%                      Custom text commands
% --------------------------------------------------------------
\newcommand{\alert}[1]{{\color{red} #1}}
\newcommand{\Alert}[1]{\textbf{\alert{#1}}}

\newcommand{\hl}[1]{\textit{#1}}
\newcommand{\Hl}[1]{\textbf{\textit{#1}}}
\newcommand{\HL}[1]{\textbf{#1}}

\newcommand{\code}[1]{\texttt{#1}}



%                    JS code formatting with minted
% --------------------------------------------------------------
% \newenvironment{longlisting}{\captionsetup{type=listing}}{}

% Define line numbering style
% \renewcommand\theFancyVerbLine{\color{bluegray400} \small \arabic{FancyVerbLine}}

% \js- команда для форматування коду прямо в рядку тексту.
% \NewDocumentCommand{\js}{v}{%
%  \usemintedstyle{friendly}%
%  \mintinline{js}|#1|
%}

% Оточеня jscode для форматоування коду як лістингу
% \newenvironment{jscode}
%    {\VerbatimEnvironment
%    \usemintedstyle{friendly}
%     \begin{minted}[breaklines, obeytabs, tabsize=4, xleftmargin=24pt, linenos, bgcolor=grey50]{js}}
%    {\end{minted}}

% Команда для включення вихідних файлів програм мовою c++
% \newcommand{\inputjs}[1]{%
%     \inputminted[breaklines, tabsize=4, obeytabs, xleftmargin=24pt,linenos]{js}{#1}%
%}  

%                    Rust code formatting with minted
% --------------------------------------------------------------
\newenvironment{longlisting}{\captionsetup{type=listing}}{}

% Define line numbering style
\renewcommand\theFancyVerbLine{\color{bluegray400} \small \arabic{FancyVerbLine}}

% \rs- команда для форматування коду прямо в рядку тексту.
\NewDocumentCommand{\rs}{v}{%
    %\usemintedstyle{friendly}%
    \usemintedstyle{vs}%
    \mintinline{rust}|#1|
}

\newenvironment{cppcode}
{\hrule\vspace{4pt}
    \VerbatimEnvironment%
    %\usemintedstyle{tango}%
    \usemintedstyle{vs}%
    \fontsize{10pt}{12pt}\selectfont%
    \setlength\partopsep{-\topsep}%
    \begin{minted}[breaklines, obeytabs, tabsize=4]{cpp}}
    {\end{minted}\vspace{4pt}\hrule}


% Команда для включення вихідних файлів програм мовою C++
\newcommand{\inputcpp}[1]{{% 🛠️ Добавили открывающую скобку тут
        \usemintedstyle{vs}
        \fontsize{10pt}{12pt}\selectfont
        \inputminted[breaklines, obeytabs, tabsize=2]{cpp}{#1}%
    }} % 🛠️ И закрывающую скобку тут

%                    Java code formatting with minted
% --------------------------------------------------------------

% Команда для автоматического включения файлов Java с мелкими номерами строк и рамками
\newcommand{\inputjava}[1]{{%
    \usemintedstyle{vs} % Используем стиль VS, как в C++
    \fontsize{8pt}{10pt}\selectfont % Размер шрифта самого кода
    
    % Локальное переопределяем размер номеров строк ТОЛЬКО внутри этой команды
    \renewcommand\theFancyVerbLine{\color{bluegray400} \scriptsize \arabic{FancyVerbLine}}%
    
    \inputminted[
        %frame=lines,      % Горизонтальные рамки сверху и снизу
        %framesep=2mm,     % Отступ кода от рамок
        baselinestretch=0.9,
        linenos,          % Включаем нумерацию
        breaklines,       % Перенос длинных строк
        tabsize=4         % Табуляция для Java
    ]{java}{#1}%
}}
